\def\week{87}
\def\topic{Voting and agendas}
\pdfinfoomitdate=1
\pdftrailerid{}
\pdfsuppressptexinfo=15
\documentclass[11pt,letterpaper]{article}
\usepackage[margin=.68in,headheight=16pt,headsep=18pt,footskip=24pt]{geometry}
\usepackage[T1]{fontenc}
\usepackage{lmodern,amsmath,amssymb,tikz,fancyhdr,array}
\usepackage[hidelinks]{hyperref}
\pagestyle{fancy}\fancyhf{}
\fancyhead[L]{\small Week \week / \topic / Facilitator}
\fancyfoot[L]{\footnotesize Bellingham Math Circle / Week \week / W\week-F-v2}
\fancyfoot[R]{\footnotesize\thepage}
\renewcommand{\headrulewidth}{.4pt}
\setlength{\parindent}{0pt}\setlength{\parskip}{7pt}
\newcommand{\heading}[1]{\par\vspace{4pt}{\large\bfseries #1}\par}
\newcommand{\sub}[1]{\par\textbf{#1}\par}
\begin{document}
\heading{Mathematical overview}
Three fixed sincere rankings $ABC$, $BCA$, $CAB$ give a majority cycle: $A$ beats $B$, $B$ beats $C$, and $C$ beats $A$, each 2--1. A vote chooses the higher of the two currently offered options on each card. Every individual ranking is consistent; the pairwise majority comparisons need not form a consistent ranking.

If two options vote first and their winner meets the waiting option, this same profile lets any option win by waiting. No card or ballot rule changes. Conversely, any option that beats both others wins every such three-option knockout. This finite example does not prove Arrow's theorem and makes no claim about fairness or real elections.

\textbf{Five fixed distance voters on one common line have a strict winner at the third favorite.} Order favorites $p_1\leq\cdots\leq p_5$ with multiplicity and assume $p_3$ is available. Against any $q<p_3$, voters 3,4,5 strictly prefer $p_3$; against any $q>p_3$, voters 1,2,3 strictly prefer it. Equal-distance ties among the other voters cannot defeat these three votes. This proves a unique strict pairwise winner against every other available location. A favorite at zero is a voter. Repeated favorites are separate voters, not one position. The middle favorite need not be the mean or the midpoint of the printed line.

For four favorites 1,3,7,9 on the line, every point in $[3,7]$ never loses a comparison, but no point beats every other. For the integer board these weak winners are 3,4,5,6,7. The even-voter result is different from saying every vote ties. With three distinct candidate positions on a line, the middle candidate cannot be any distance voter's strict worst choice. Since the cyclic profile puts each candidate last once, no ordering of candidate positions realizes all three cards by nearest-distance preferences.

\sub{Discovery and readiness}
Approximate grades 2--3 can retain imaginary cards while physically voting and changing the waiting option. Upper readiness adds profile design and explanation. The separate nearest-location visit usually needs grades 3--5 readiness: compare distances, retain fixed favorites, distinguish proposal markers and account for ties. Repeated favorites, even voters and impossible preference models are return investigations, not compulsory first-hour work.

Experiments are completed votes and tested agendas. The three agenda choices exhaust the first profile's possibilities. Testing a few favorite arrangements only suggests the median rule; the three-voter coalition on either side proves it generally under the assumptions above.

\sub{K--1 route (adult-led)}
There is no independent K--1 packet. A K--1 table can do shared Problems 1--2 with an adult reading the cards; children need only recognize A, B and C and count to three. Use the fixed cards, option tiles and vote tokens from materials p.~1 and the tally and knockout mats from p.~4. For each pair the adult reads one card aloud (``Voter 1 likes A best, then B, then C. A or B?''), a child puts that voter's token in the chosen well, and the children count. Then the children choose who waits and move the tiles on the knockout mat while the adult reads each vote. Plain-words key: A beats B, B beats C and C beats A, each two tokens to one, so no letter beats both others. Whichever letter waits wins: B and C first makes A win, C and A first makes B win, A and B first makes C win. Leave out profile design and the line visit. This route is untested.

\sub{Status and limits}
Prepared materials are unpiloted. Cutting, stacking markers, resetting a knockout, tie handling and timing have not been physically rehearsed. The model assumes sincere choices in each current pair, fixed preferences, a common line and available median option. It does not model strategic anticipation or arbitrary human preferences.
\newpage
\heading{Preparation and launch}
Prepare three older-table pair kits if both tables use the activity. The upper table rotates the third child as referee. Nobody needs to impersonate five voters from memory: cards or favorite markers remain visible and each voter has one token. The anchored adult may read or record while the child chooses each ballot. Repeated adult selection of votes means the route is not yet independent.

For the ranking visit use the supplied three ranking cards, option tiles, one token per voter, tally mat and knockout mat. For profile design add six blank erasable ranking cards. For the line visit add the 0--10 working board, five distinct favorite markers, two differently shaped proposal markers X/Y, separate X/Y tally labels and V1--V5 vote tokens. All five tokens must fit in any tally well, including Tie. Repeated favorites use separate parking lanes. The four-voter visit removes both V5's favorite and its vote token. Print only the materials needed that day.

Budget 15--25 minutes for one ranking visit with existing counters/erasable cards, or about 35--50 minutes to cut all three complete kits; add 10 minutes to rehearse an unavailable favorite, a knockout reset and a tied line ballot. These are untested estimates. Print US Letter at actual size and verify the line spacing and token fit. Favorite markers stay fixed while proposal markers and vote tokens move. The 170 by 90 mm tally mat has three 54 by 54 mm wells. Its matching 18 mm targets have 2 mm boundary clearance and at least 4.6 mm between placed tokens; the separate tally labels are 23 mm squares. These are digital dimensions, not a physical fit test.

\sub{First launch: fixed cards}
After movement and a minute handling tiles, offer two options against one non-task ranking card. Ask which offered option is higher; physically put that voter's token beside it. Let children do the other cards, including a card whose first choice is absent. Finish the whole vote on the same unchanged three cards. Explain that these are imaginary voters, not the children's own preferences. Later show first-round winner and waiting option advancing into one completed final vote before asking children to choose agendas.

\sub{New launch for the line visit}
Use the printed non-task favorites 0,1,5,8,10 and simultaneous proposals X=2, Y=8. Keep that arrangement fixed throughout. The X and Y proposal markers stay on the board; the separate X/Y tally labels name the wells. For each voter compare distances to both proposals, then put its token in X, Y or Tie. A tied voter has one token in Tie, never one for each proposal. V1 and V2 choose X, V3 ties, and V4 and V5 choose Y: the result is 2--2 with one voter in Tie. Complete the tally before changing any proposal; favorites remain on their parking lanes. Give children time to move proposals and test votes before introducing any general middle-favorite claim.

\sub{Pacing}
Treat agenda and line investigations as separate visits. A first agenda visit can use 10--15 minutes on pair votes, 15--20 on different waiting choices, then optional profile design. Upper Problem 1 continues the same visit after shared Problem 2: it needs only the card rules, and its first set shows a letter that beats both others before shared Problem 3 asks children to build one. A line visit can spend most table time testing a fixed favorite arrangement before inventing others. Even-voter and impossibility questions are later choices. Record a chosen agenda or final favorite/tally snapshot, not simultaneous ranking transcription and distance tables.
\newpage
\heading{Shared packet: fixed cards and agendas}
\sub{Problem 1: three pairwise votes}
With cards $ABC$, $BCA$, $CAB$, V1 and V3 choose A over B; V1 and V2 choose B over C; V2 and V3 choose C over A. Every result is 2--1. Keep all three cards visible, including the unused option in each ranking. A voter's favorite being absent does not remove that voter: compare the two options that are actually offered. There is no overall top option in these pair results, even though each individual card has a strict first, second and third choice.

The non-task worked example uses $DFE$, $EFD$, $FED$. When D and E are offered, V1 chooses D and V2/V3 choose E, so E wins 2--1. The fixed cards remain unchanged in the next demonstration: E then meets waiting F; V1/V3 choose F and V2 chooses E, so F wins 2--1. A child who votes for an unavailable F in the first round needs the available pair made explicit, not a new preference card.

\sub{Problem 2: choose who waits}
Keep the Problem 1 cards. To make A win, first compare B with C; B wins, then A beats B. To make B win, first compare C with A; C wins, then B beats C. To make C win, first compare A with B; A wins, then C beats A. Thus the waiting option wins each time. These three choices exhaust all unordered first pairs. Reversing the left/right position of a pair changes no ballot under the rule.

Distinguish the first-round loser, first-round winner, waiting option and final winner. The knockout mat retains those states. Reset the movable options between attempts; do not reset or rewrite the preference cards. A child can record just the first pair and final winner. This is an agenda effect under sincere pairwise voting, not evidence that voters changed their minds or voted strategically.

\sub{Problem 3: design two profiles}
A suitable three-distinct-card profile with one winner for every agenda is $ABC$, $ACB$, $BAC$. A beats B 2--1 and C 3--0. If A waits, it beats the first winner; if A plays first, it wins both of its comparisons. Three different cards cannot share a first choice, since only two rankings start with each letter. The cyclic profile of Problem 1 is already a valid three-distinct-card agenda-dependent example; children may reproduce or relabel it. Every working set uses exactly the cards $ABC$, $BCA$, $CAB$ or $ACB$, $CBA$, $BAC$, in some order. Each card agrees with at most two of the cycle's three winning pairs, and each winning pair needs two cards, so every card agrees with exactly two. ``Is there a different one?'' makes a good return question.

Do not require every child to enumerate all possible profiles. Accept a construction with its checkable pair votes and agenda argument. For a held hint, ask what an option would need to do against each of the other two. Blank ranking cards and movable slips preserve each invented profile without repeatedly transcribing it. If two identical cards appear, remind the child of this task's distinct-card condition; a repeated ranking would still be a legitimate voter in the general model.
\newpage
\heading{Line packet: a winner among locations}
All line questions use fixed voters, two simultaneous proposal locations, absolute distance, and a separate Tie destination for equal-distance voters. Vote tokens move; favorite markers do not. The printed 0--10 ticks are available choices, not a demand to list all 55 pairs. This packet starts a new visit, not an extra recording layer during the card investigations.

\sub{Problem 1: favorites 0,2,4,7,10}
Location 4 beats every other integer location. Against 0, the voters at 4,7,10 choose 4; the voter at 0 chooses 0; the voter at 2 ties. The count is 3--1 with one Tie. Against 10 the voters at 0,2,4 choose 4; the voter at 10 chooses 10; the voter at 7 ties: again 3--1 with one Tie. Against each of 1,2,3,5,6,7,8,9, the count is 3--2. A tied individual ballot is not an overall tied result.

The middle favorite is the answer, but do not state that before children compare locations. If a group gets bogged down testing every pair, invite them to choose a promising proposal and let a partner try to defeat it. Save the side-coalition explanation for a discovered conjecture or Problem 3. Distinguish a successful challenge from a tied vote: a strict winner must win every comparison, so either result refutes that claim.

\sub{Problem 2: repeated and uneven favorites}
For favorites 0,0,0,8,10, the winner is 0. Three different voters at zero each prefer zero to every distinct proposal. Their markers occupy different lanes and their votes remain separate. For 1,4,4,4,9, the winner is 4 for the same three-voter reason. For 0,1,6,9,10, the winner is 6. This last configuration distinguishes the middle favorite from the board midpoint 5 and arithmetic mean $26/5$. Its voters are numbered out of left-to-right order (V4 is at 6), so no voter number points to the winner.

Do not combine coincident favorites into one token or discard a voter at zero. These examples deliberately retain multiplicity. A held adult prompt is to point to each voter's identity and ask which one has not yet voted, without selecting their vote. A tally that contains fewer than five tokens is unfinished, even if the outcome is already forced.

\sub{Problem 3: five favorites anywhere on the ticks}
Sort the five favorites with repeats retained: $p_1\leq p_2\leq p_3\leq p_4\leq p_5$. The third, $p_3$, wins against every distinct proposal q. If q is left of $p_3$, each of voters 3,4,5 is at or right of $p_3$, so moving farther left to q strictly increases their distance. Those three vote for $p_3$. If q is right of $p_3$, voters 1,2,3 supply the analogous three votes. The remaining two votes, including any ties, cannot reverse the majority.

This proof works on a continuous line too, provided the median location is an available option. It does not prove a result for arbitrary rankings, a circle, two-dimensional distance, or a restricted menu excluding the median. Testing every five-favorite multiset on these eleven ticks checks 3,003 arrangements but is not the reason the general theorem holds. Children can demonstrate the two defending groups using favorite markers and a proposal on either side, without algebra. If a child needs a direction, a held prompt is: ``Does the third favorite from the left always beat every different location? Show why, or find a set where it fails.''
\newpage
\heading{Return investigations and the agenda continuation}
\sub{Line Problem 4: four favorites 1,3,7,9}
Exactly the tick locations 3,4,5,6,7 never lose. No location strictly beats every other location. To see the first assertion, take x between 3 and 7. Against any proposal left of x, voters at 7 and 9 favor x, so x has at least two votes and cannot lose among four voters. Against any proposal right of x, voters at 1 and 3 defend x in the same way. Two distinct proposals both inside $[3,7]$ split the four votes 2--2. Hence none of those locations strictly beats all others.

A tick below 3 loses to 3: the voters at 3,7,9 all prefer 3. A tick above 7 similarly loses to 7, supported by voters at 1,3,7. This excludes all outside ticks and proves completeness. The same argument gives the whole closed interval $[3,7]$ if all real locations are available. There are still winning comparisons elsewhere; ``no universal strict winner'' does not mean every comparison ties.

Remove both the fifth favorite marker and its vote token. Changing only one creates a different, inconsistent procedure. Let children preserve several tied comparisons before asking for an explanation of all the weak winners.

\sub{Upper Problem 1: can the agenda change the winner?}
Use it in the first agenda visit, after shared Problem 2; it needs only the card rules. The first profile is $ABC$, $ACB$, $BAC$. A beats B 2--1 and C 3--0, so A wins every agenda by the argument in shared Problem 3. The other pair has B beating C 2--1, but its result cannot remove A.

The second profile is $ACB$, $CBA$, $BAC$. B beats A, A beats C, and C beats B, each 2--1. Thus each waiting option can win: first A/C then waiting B; first B/C then waiting A; first A/B then waiting C. The task asks children to use pair votes to explain the possible agenda effect, not merely replay every knockout. In the three-option, odd-voter strict-ranking model, the pair comparisons form either a transitive order with a universal winner or a three-cycle. This observation should not be silently extended to larger menus, where more patterns occur.

\sub{Upper Problem 2: can a common line produce the cycle?}
No. Put the three distinct candidate locations in order $\ell<m<r$. For any voter favorite p at or left of m, the rightmost candidate r is strictly farther away than m. For any p at or right of m, the leftmost candidate $\ell$ is strictly farther away than m. At least one extreme is therefore farther from p than the middle candidate; m cannot be that voter's strict worst choice. This holds for arbitrary real positions and favorites, including favorites outside the candidate interval.

The cards $ABC$, $BCA$, $CAB$ each put a different candidate last. Whichever candidate occupies the middle would have to be strictly last for one of the voters, contradicting the preceding fact. Coincident candidates or tied distances cannot rescue a strict ranking. The worked conversion on the student page shows how one fixed favorite produces a ranking; failed experimental placements alone would not prove impossibility. This gives a concrete limit of the nearest-location model, not a judgment about people's preferences.
\newpage
\heading{Print map, provenance and use record}
The shared packet has three problems on three pages; the separate line packet has four on four pages; the upper packet has two on two pages. Shared and upper profile questions overlap deliberately so each packet can stand alone. Guide labels identify the packet before the problem number. Do not hand every page to every pair at once.

The five materials pages supply: (1) the cyclic ranking cards, option tiles and vote tokens; (2--3) two sets of three blank ranking cards and letter slips; (4) tally and knockout mats; (5) a five-lane favorite board and markers. Use the revised tally labels to retain X/Y proposals on the board while the wells remain identified. Check that all five vote tokens fit in any one destination with room to place and read them. Repeated favorite positions stay separated by voter lane. Print at actual size and check the scale reference; do not substitute a reduced thumbnail for a working board.

\sub{Research and adaptation}
Duncan Black, \emph{On the Rationale of Group Decision-making}, Journal of Political Economy 56(1), 1948, pp.23--34, especially section II, pp.26--28, is the primary source for the one-dimensional single-peaked majority result. The production research read those pages in an original scan. This packet uses the narrower symmetric distance-preference model, whose odd-voter conclusion is proved here directly. It does not assert that all single-peaked preferences are distance rankings or import broader conclusions without their hypotheses.

Daron Acemoglu, \emph{14.770 Introduction to Political Economy, Lectures 1--2}, MIT OpenCourseWare, slides 20--23, provides a modern statement of the median-voter argument and assumptions. \href{https://ocw.mit.edu/courses/14-770-introduction-to-political-economy-fall-2017/bdc7ef5c8617eaf573dc28da5c2e19ef_MIT14_770F17_lec1_2_acemoglu.pdf}{MIT course resource}. The majority-cycle and agenda constructions are locally chosen finite examples, checked independently. They are not an elementary proof of Arrow's theorem. The original planning bibliography's Stodder resource was not freshly available during this production pass; no new reading of it is claimed.

Givental, Nemirovskaya and Zakharevich, \emph{Math Circle by the Bay}, Preface viii--x, supports preserving deep revisitable themes, independent work and concrete materials. Our exact route split, staffing estimates and kit design are editorial decisions, not reported classroom outcomes. Earlier median work concerns sums of travel distances; this week asks pairwise majority questions. Weeks 92--97 do not duplicate this agenda/line-model sequence. No borrowed illustrations, articles or downloaded references are bundled with the editable sources.

\sub{What to observe next}
Record the packet/problem used, whether each voter remained represented, whether children distinguished a personal tie from an aggregate tie, and an explanation in a child's words. Note separately any adult-selected ballots or repeated reset help. These are observations; possible reasons for confusion are hypotheses until checked. Retain one invented profile or favorite/tally snapshot for a later visit. Cutting, token fit, reset procedures, pacing and classroom use remain physically untested and unpiloted. Git publication does not imply a Drive upload or a synchronized classroom copy.
\end{document}
